\section{Conclusion}\label{sec:conclusion}

We took a census of \kvcCensusRows{} open-weight checkpoints from
\kvcOrgCount{} organizations, reduced to \kvcDedupRows{} family-level groups
(\kvcGeoRows{} distinct attention geometries), accounted their KV caches from
configurations alone, and validated the accounting against reference
implementations. Four verdicts.
For RQ1, per-parameter KV cost is falling: the pooled 2025--2026 median is
\kvcHOneRatio{} of the 2023 median on the pre-registered unit and
\kvcHOneRatioGeo{} under geometry deduplication (H1 supported on both), though
not monotonically under the family unit, because the 2025 release mix and
re-released same-geometry fine-tunes moved against the trend; excluding
multimodal checkpoints flips the verdict (\kvcHOneRatioNoMM), so the finding is
real but fragile.
For RQ2, parameter count is a weak predictor of KV footprint ($R^2=$
\kvcParRsq) while attention geometry alone explains \kvcGeomRsq, and MoE models
carry \kvcHThreeRatio{} times the dense per-parameter KV cost (H2 and H3
supported).
For RQ3, the KV-parity context grew \kvcHFourRatio{} under the pre-registered unit
and \kvcHFourRatioGeo{} under geometry deduplication (H4 refuted on the first,
supported on the second): KV savings and weight growth offset each other, and
release mix sets the median, even as KV-dominance at 32k context receded.
For RQ4, the analytic accounting matches reference implementations on
\kvcValWithin{} of \kvcValRows{} validation rows within 1\%, across
\kvcValClasses{} attention classes.

The practical rule is the title: size serving memory by attention geometry, not
by parameter count. Read the layer count, KV heads, head dimension, window
sizes, and state shapes from the configuration; compute the cache; then pick the
quantization levers the geometry leaves room for. On a fixed unified-memory
device this computation, not the parameter count on the model card, decides how
many long contexts the hardware can hold.

Three extensions follow. Measure real runtimes at scale, adding the allocator
and kernel effects the analytic model omits. Extend to closed models where
configuration metadata becomes available. Re-run the census periodically:
window share, latent compression, and hybrid recurrence are still moving
quickly, and the 2026 medians already rest on a partial year.
